% HD-F1-T19. GPT-6 Astra Pro, /root/transport; 22 November 2026 UTC.
\section[HD-F1-T19: defining fields at most 27]{HD-F1-T19: actual rank-two subgroup maps remove defining field sizes 29, 31 and 32}
\label{sec:f1-rank-two-subgroup-field-exclusion}

Every remaining T2-OTHER source with almost simple projective
coefficient image of Lie type has defining field size
\(q\le27\). In addition to T18, the proof removes exactly
\(q=29,31,32\). It uses explicit restriction maps from
rank-two subgroups to the original coefficient representation.
The degree lower bounds needed for those restrictions are
proved below by finite character projections, without an
input from the classification of minimum representation degrees.

\paragraph{1. Original central-cover map and subgroup inputs.}
For each of these three field sizes, retain the surjection
\(\beta:\widetilde S\to L\subset G\) from T16--T18.
Its kernel is central, and the centre of \(\widetilde S\)
has order prime to the defining characteristic \(p\).
These facts follow from the already checked generic-cover
scope in Malle--Testerman, Theorem~24.17,
Remark~24.19 and Table~24.3; every exceptional defining
field size is at most nine.

The actual auxiliary maps below have the form
\[
 H\xrightarrow{\alpha}\widetilde S
    \xrightarrow{\beta}L\hookrightarrow G
    \hookrightarrow GL(V_{\mathbf C}),
 \qquad H=SL_3(q)\ \text{or}\ SU_3(q).
 \tag{F1T19.1}
\]
For root-subgroup constructions, \(\alpha\) is the
central isogeny from the simply connected rank-two
subsystem onto its image. For block or listed subgroup
constructions it is the inclusion. In every case
\(\alpha\) is injective on the indicated unipotent
subgroup, and its nonidentity unipotent elements remain
noncentral in \(\widetilde S\). Thus \(\beta\alpha\)
does not kill any of those elements. In particular the
representation in (F1T19.1) is nontrivial on the subgroups
used in paragraphs~2 and~3.

The root-subgroup input is Malle--Testerman,
Theorem~13.6 and Corollary~13.7, pp.~106--107:
a closed root subsystem gives its actual semisimple
subgroup generated by the corresponding root groups.
The complete construction proof was read. Its simply
connected covering map is Proposition~9.15, p.~71;
here the characteristic is prime to the \(A_2\)
fundamental-group order three, so the supplied proof
of that proposition applies directly. The pinned graph
maps are Theorem~11.12, pp.~89--90. For a closed
subsystem \(A_2\) with simple roots \(\alpha_0,\beta_0\),
the simply connected covering morphism is specified on
root generators by
\[
 x_{12}(c)\longmapsto x_{\alpha_0}(c),\quad
 x_{23}(c)\longmapsto x_{\beta_0}(c),\quad
 x_{13}(c)\longmapsto x_{\alpha_0+\beta_0}(\epsilon c),
 \qquad \epsilon\in\{1,-1\},
 \tag{F1T19.2}
\]
and by the compatible negative-root generators. The sign
is chosen to agree with the Chevalley commutator relation.
The covering morphism has central kernel; when these
root parameters are fixed by the graph part of Frobenius,
(F1T19.2) commutes with the \(q\)-power Frobenius and
restricts to the finite-group map in (F1T19.1).

\paragraph{2. An actual \(SL_3(q)\) restriction costs at least \(q^2-1\).}
Inside \(SL_3(q)\), take
\[
 U=\left\{\begin{pmatrix}I_2&v\\0&1\end{pmatrix}:
                              v\in\mathbf F_q^2\right\},
 \qquad
 D=\left\{\begin{pmatrix}A&0\\0&1\end{pmatrix}:
                              A\in SL_2(q)\right\}.
 \tag{F1T19.3}
\]
The exact conjugation action of \(D\) on \(U\) is
\(v\mapsto Av\). For each additive character \(\eta\)
of \(U\), the operator
\[
             e_\eta=\frac1{|U|}
                   \sum_{u\in U}\eta(u)^{-1}\rho(u)
 \tag{F1T19.4}
\]
is the projection onto the actual \(\eta\)-weight space
of \(V_{\mathbf C}\). Its conjugate by \(\rho(A)\)
is the projection for the corresponding dual character.
Using the nondegenerate field-trace pairing to identify
the character group of \(\mathbf F_q^2\) with its dual
vector space, \(SL_2(q)\) is transitive on all \(q^2-1\)
nonzero characters. Nontriviality on \(U\) makes at
least one such weight space nonzero; conjugation makes
all of them nonzero and of equal dimension. Therefore
the original coefficient dimension satisfies
\[
                              m\ge q^2-1.
 \tag{F1T19.5}
\]
No irreducibility of the restriction is needed.

\paragraph{3. An actual \(SU_3(q)\) restriction costs at least \(q(q-1)\).}
Use the Hermitian form with antidiagonal matrix
\(J=\left(\begin{smallmatrix}0&0&1\\0&1&0\\1&0&0\end{smallmatrix}\right)\)
on \(\mathbf F_{q^2}^3\). The following matrices belong
to its special unitary group:
\[
 u(a,b)=\begin{pmatrix}1&a&b\\0&1&-a^q\\0&0&1\end{pmatrix},
 \quad b+b^q=-a^{q+1},
 \qquad
 t(z)=\operatorname{diag}(z,z^{q-1},z^{-q}),
 \quad z\in\mathbf F_{q^2}^{\times}.
 \tag{F1T19.6}
\]
The \(u(a,b)\) form a group \(U\) of order \(q^3\),
with centre
\[
 Z_U=\{u(0,b):b+b^q=0\},\qquad |Z_U|=q,
 \tag{F1T19.7}
\]
as follows directly from the formulas
\[
 \begin{split}
 u(a,b)u(c,d)&=u(a+c,b+d-ac^q),\\
 [u(a,b),u(c,d)]&=u(0,ca^q-ac^q).
 \end{split}
 \tag{F1T19.8}
\]
For every \(a\ne0\), the map
\(c\mapsto ca^q-ac^q\) is an onto \(\mathbf F_q\)-linear
map to the trace-zero line, with kernel \(\mathbf F_q a\).

Let \(\eta\ne1\) be an additive character of \(Z_U\).
Every irreducible \(U\)-constituent with central
character \(\eta\) has dimension exactly \(q\).
Indeed, for any noncentral \(u(a,b)\), surjectivity
in (F1T19.8) gives a conjugating element that changes
it by a central element on which \(\eta\ne1\).
Its character consequently vanishes off \(Z_U\).
If its dimension is \(d\), the character on the centre
is \(d\eta\); the averaged Hom projection for that
irreducible constituent has trace one, and gives
\[
                 1=\frac1{q^3}\sum_{v\in U}|\chi(v)|^2
                    =\frac{qd^2}{q^3},
                 \qquad d=q.
 \tag{F1T19.9}
\]
This calculation allows nonfaithful additive characters
of \(Z_U\) when \(q\) is not prime: the commutator map
is onto the whole trace-zero line before applying
\(\eta\).

The exact torus conjugation on the centre is
\[
          t(z)u(0,b)t(z)^{-1}=u(0,z^{q+1}b).
 \tag{F1T19.10}
\]
The norm \(z\mapsto z^{q+1}\) is onto
\(\mathbf F_q^\times\). Its action is transitive on
the \(q-1\) nontrivial additive characters of \(Z_U\).
The representation (F1T19.1) is nontrivial on \(Z_U\),
so at least one of their projected weight spaces is
nonzero. All \(q-1\) are then nonzero and have equal
dimension, each at least \(q\) by (F1T19.9).
It follows that
\[
                              m\ge q(q-1).
 \tag{F1T19.11}
\]
Again this is a statement about the full actual
restriction, with no irreducibility requirement.

\paragraph{4. The maps cover every remaining family at these three field sizes.}
Use the usual pinned root numbering, as in
Malle--Testerman, Table~13.1. The following list specifies
the subgroup maps, rather than substituting a natural
module for the original coefficient space.
\begin{itemize}
\item For the split types \(A_r\) with \(r\ge2\),
\(B_r,C_r\) with \(r\ge3\), and \(D_r\) with
\(r\ge4\), take the closed \(A_2\) subsystem on the
first two adjacent equal-length simple roots. For the
split types \(F_4,E_6,E_7,E_8\), take adjacent
equal-length simple roots; in \(F_4\) one can take
the two adjacent long simple roots. Formula (F1T19.2)
then gives the \(SL_3(q)\) map.
\item For \({}^2A_r(q)=SU_{r+1}(q)\), \(r\ge2\),
choose a nondegenerate Hermitian three-dimensional
summand of the defining Hermitian space and let
\(SU_3(q)\) act on it and trivially on its orthogonal
complement. This is an actual block inclusion; the
determinant remains one.
\item For \({}^2D_r(q)\), \(r\ge4\), the graph
automorphism swaps the two end nodes and fixes the
first two adjacent simple roots. For \({}^2E_6(q)\),
it fixes the adjacent simple roots numbered 2 and 4.
Their pinned root parameters are fixed by the graph
part, so (F1T19.2) gives \(SL_3(q)\) in both cases.
\item For \({}^3D_4(q)\), let the central simple root
be \(\alpha_2\), and let triality cycle the other three
simple roots. The two roots
\[
       \alpha_2,\qquad
       \gamma=\alpha_1+\alpha_2+\alpha_3+\alpha_4
 \tag{F1T19.12}
\]
are fixed by triality and generate a closed \(A_2\)
subsystem. For an explicit check, in the standard
\(D_4\) coordinates they are \(e_2-e_3\) and
\(e_1+e_3\), with sum \(e_1+e_2\). They have
equal squared length two and inner product minus one.
A pinned graph automorphism acts on a fixed Chevalley
root parameter by a sign; its order three forces this
sign to be one. Thus (F1T19.2) commutes with the
triality-twisted Frobenius and gives \(SL_3(q)\).
\item For \({}^2F_4(q)\), only \(q=32\) among the
three field sizes is possible. Use the subgroup
\(SU_3(q)\) inside \(SU_3(q):2\) explicitly listed
in Wilson, \emph{On the simple groups of Suzuki and
Ree}, \emph{Proc. London Math. Soc.}~107 (2013),
680--712, Section~9.4, Theorem~20(iii), author PDF
p.~29, citing Malle's subgroup theorem. Only this
subgroup inclusion is used, not exhaustiveness of the
maximal-subgroup list. Here the simply connected
fixed-point group has trivial centre, so it maps
isomorphically to the actual perfect layer \(L\).
\end{itemize}
The only Lie-type families not in this list are
\(A_1\), \(B_2=C_2\), \(G_2\), \({}^2B_2\) and
\({}^2G_2\), already excluded at these parameters by
T14. Low-rank identifications of types \(D\) and
\({}^2D\) are the displayed linear or unitary cases.
Thus the list covers every remaining simple socle
at \(q=29,31,32\).

\paragraph{5. Original-rank contradiction and retained framework maps.}
Both lower bounds (F1T19.5) and (F1T19.11) are at
least \(29\cdot28=812\). This exceeds all three
original-rank ceilings \(87,46,38\) in T18. Hence
none of the three field sizes is possible. Combining
this with \(q\le32\), and observing that the next
smaller prime power is 27, proves \(q\le27\).

Every restriction in this proof is the composite
(F1T19.1) into the original \(V_{\mathbf C}\).
The character projectors (F1T19.4) and their central
analogues, the conjugation maps and the norm map
(F1T19.10) are the actual dimension-comparison maps.
No subgroup quotient curve or smaller representation
is proposed as a detector. The cover, middle-extension
sheaf, rational primitive Jacobian projector, actual
sixfold corner and restricted theta polarization are
retained by identity maps. The same identity maps
preserve their Poincare biextension, parabolic pairing,
variational de Rham, Grothendieck and Adams realizations,
with polarization multiplier one. This is a further
source exclusion on the original natural-curve gate.
